\section{Future Directions}

This study opens several promising avenues for future research:

\paragraph{Exploration of Workflow Vulnerabilities and Defense Mechanisms: }Investigating how negotiation workflows can be exploited through methods such as deception is crucial, particularly in contexts like the Deal-No-Deal Game. Understanding the potential vulnerabilities within these workflows will enable the development of robust defense strategies to mitigate exploitation risks, thereby enhancing the reliability and security of negotiation protocols.

\paragraph{Strategizing in Multi-Stage Games:} Extending the analysis to multi-stage games introduces additional complexity, as the Nash Equilibrium may differ significantly from that in single-stage games. Future work should focus on how LLMs can effectively strategize over multiple stages, accounting for the dynamic evolution of the game state and the opponent's actions. This includes developing algorithms that can anticipate future moves and adjust strategies accordingly to maintain rationality and optimize outcomes.

\paragraph{Development of Meta-Strategy:} LLMs should be equipped with the capability to determine when to employ a particular workflow and when to adapt or forgo it. This necessitates the creation of a meta-strategy or meta-workflow that guides the selection of appropriate negotiation strategies based on the specific context, the agent's own abilities, the stage of the game, and the behavior of the opponent. Implementing such a meta-strategy would enhance the adaptability and effectiveness of LLMs across diverse negotiation scenarios.

\paragraph{Alignment with Agent Interests and Stance Adoption:} Currently, LLMs are generally aligned to be helpful and honest, lacking a personalized stance or specific interests. To function as proficient negotiation agents, it is imperative for LLMs to learn and adopt stances that reflect the interests they represent. This involves training LLMs to understand and advocate for particular objectives, thereby balancing their general alignment with the capacity to pursue defined goals within negotiations. Developing methods to instill a clear understanding of their own interests will enhance the LLMs' ability to engage in strategic decision-making and achieve desired outcomes.

These future directions aim to refine the strategic reasoning and negotiation capabilities of LLMs, ensuring they can operate effectively and rationally in complex, real-world scenarios. By addressing these areas, we can advance the development of LLMs as robust agents capable of navigating intricate strategic environments while safeguarding against potential vulnerabilities.
